\documentclass[11pt,a4paper]{article}
\usepackage[margin=25mm]{geometry}
\usepackage{amsmath,amssymb,amsthm,hyperref}
\newtheorem{lemma}{Lemma}
\newcommand{\gen}[1]{\langle#1\rangle}
\title{Fixed prefixes at a unique orbit crossing}
\author{Human Compression project: L044 working continuation}
\date{}
\begin{document}\maketitle
Let $G=\gen{\rho_0,\ldots,\rho_{r-1}}$ be a faithful transitive
string action by involutions. Suppose colour $f$ has a unique edge
$\{a,b\}$ whose endpoints are in different orbits of
$G_f=\gen{\rho_j:j\ne f}$. The single-crossing lemma in
\emph{Paired orbit cuts in string actions} gives exactly two
$G_f$-orbits $A,B$, containing $a,b$, and says that every
$\rho_j$ with $|j-f|>1$ fixes both $a,b$. Write
$G_{>l}=\gen{\rho_j:j>l}$ and define $G_{<l}$ dually.

\begin{lemma}[A fixed prefix to an active colour]\label{prefix}
Suppose $l>f$ and $\rho_l$ acts nontrivially on $A$. There is a
path in the full permutation graph from $a$ to a point moved by
$\rho_l$, using only colours less than $l$ and different from $f$.
Every vertex of this path is fixed by $G_{>l}$.
The dual assertion for $l<f$ uses only colours greater than $l$
and different from $f$, and its vertices are fixed by $G_{<l}$.
The same assertions hold with $b,B$ in place of $a,A$.
\end{lemma}
\begin{proof}
Put $H=\gen{\rho_j:j<l,\ j\ne f}$ and $R=G_{>l}$.
Every generator of $R$ fixes $a$, because $j>l>f$ implies
$j-f\ge2$. Every generator of $R$ also commutes with every
generator of $H$. Consequently $R$ fixes the whole orbit $a^H$
pointwise. If $\rho_l$ fixed that orbit pointwise as well, then
the orbit would be invariant under all the generators of $G_f$.
It would therefore equal $A=a^{G_f}$, contradicting the assumed
nontrivial action of $\rho_l$ on $A$. Some $H$-word thus takes
$a$ to a point moved by $\rho_l$. Its successive vertices lie
in $a^H$ and are fixed by $R$. Erasing loops gives the required
simple path. Reversing the index order proves the dual claim;
the argument at $b$ is identical.
\end{proof}

\begin{lemma}[Consecutive-colour lower bound]\label{length}
Under the hypotheses of Lemma~\ref{prefix} with $l>f$, every
path from $a$ to a point moved by $\rho_l$ that avoids colour $f$
contains each of $f+1,\ldots,l-1$. In particular appending the
final $l$-edge to a shortest prefix in Lemma~\ref{prefix} gives
a simple path with at least $l-f$ edges. Its prefix, excluding
that final edge, is pointwise fixed by $G_{>l}$.
\end{lemma}
\begin{proof}
If a colour $k$ with $f<k<l$ were missing, factor the path word
into a word in generators of index less than $k$ and a word in
generators of index greater than $k$. The two sets commute.
Every generator of index greater than $k$ fixes $a$, because
its index is at least $f+2$. Thus the endpoint is reached from
$a$ by the lower-index word alone. The generator $\rho_l$
commutes with that word and fixes $a$, so fixes the endpoint,
a contradiction. This proves the colour and length assertions.
Choose a shortest path, among the permitted prefixes of
Lemma~\ref{prefix}, to any point moved by $\rho_l$. Every
earlier vertex is then fixed by $\rho_l$. The other endpoint
of the final $l$-edge is also moved by $\rho_l$, so cannot be
an earlier prefix vertex. The extended path is simple.
\end{proof}

\begin{lemma}[The endpoint case of source Proposition 5.19]\label{endpoint}
Suppose $f=0$, every generator is even, $r\ge4$, and
$\rho_{r-1}$ acts nontrivially on both $G_0$-orbits. Then
$n\ge2r+1$, equivalently $r\le\lfloor(n-1)/2\rfloor$.
\end{lemma}
\begin{proof}
Apply Lemma~\ref{length} in each orbit with $l=r-1$. The two
disjoint simple paths each have at least $r-1$ edges, so together
they have at least $2r$ vertices. It remains to exclude $n=2r$.

At equality the paths exhaust their respective orbits. Each has
exactly one edge of each colour $1,\ldots,r-1$, in that order:
the no-gap assertion applied to each initial segment forces the
first occurrence of every colour to follow all smaller positive
colours. Write the paths
\[
 a_0=a\mathrel{-_1}a_1\mathrel{-_2}\cdots
 \mathrel{-_{r-1}}a_{r-1},\qquad
 b_0=b\mathrel{-_1}b_1\mathrel{-_2}\cdots
 \mathrel{-_{r-1}}b_{r-1}.
\]
Because $a_0b_0$ is the unique cross-orbit $0$-edge,
$\rho_0$ preserves each tail $A'=\{a_1,\ldots,a_{r-1}\}$ and
$B'=\{b_1,\ldots,b_{r-1}\}$. The group
$H=\gen{\rho_2,\ldots,\rho_{r-1}}$ fixes $a_0,b_0$ and
preserves both tails; its action on each tail is transitive,
since the displayed edges form a path through that tail.

The subgroup $K=\gen{\rho_3,\ldots,\rho_{r-1}}$ fixes $a_1$:
each of its generators fixes $a_0$ and commutes with $\rho_1$.
Every other vertex of $A'$ is moved by one of the displayed
edges of colour at least three. This uses $r\ge4$, and shows
that $a_1$ is the unique $K$-fixed point in $A'$. The restriction
of $\rho_0$ centralises $H$, hence $K$, so must fix $a_1$.
Centralising the transitive $H$-action then forces it to fix
every point of $A'$. The identical argument fixes all of $B'$.
Therefore $\rho_0=(a_0\,b_0)$ is a single transposition,
contrary to evenness. This excludes equality.
\end{proof}

For the main L044 range $n\ge12$, ranks $r\le3$ already satisfy
the required bound, so Lemma~\ref{endpoint} covers the whole
endpoint/top-colour case. The argument does not assume full IP,
does not use the unresolved square-rigidity assertion, and does
not assert that all other source branches are complete.

\paragraph{Exact relation to the source.}
The no-gap idea comes from Cameron--Fernandes--Leemans--Mixer,
\emph{Highest rank of a polytope for $A_n$},
arXiv:1605.09173v1, Proposition~5.18, physical pages29--30.
The statement above specifies the initial orbit, the active target
colour, the side of the string which fixes the path, and the
exclusion of its final edge. The source's final sentence instead
prints $G_{>l}$ when $l<f$ and $G_{<l}$ when $l>f$ and does not
make the prefix boundary explicit. We do not identify that literal
sentence with Lemma~\ref{prefix}. No bound for the whole original,
new conjecture solution or formal verification is asserted here.
\end{document}
